% Part:sets-functions-relations
% Chapter: size-of-sets
% Section: non-enumerability

\documentclass[../../../include/open-logic-section]{subfiles}

\begin{document}

\olfileid{sfr}{siz}{nen}

\olsection{\printtoken{S}{nonenumerable} سیٹ}

\begin{editorial}
  ایہ حصہ \olref[enm]{sec} وچ دتی تعریف نوں ورت کے $\Bin^\omega$ تے
  $\Pow{\PosInt}$ دے گنتی وار فہرست نہ بنن جوگ ہون نوں ثابت کردا اے۔
  اینوں \olref[enm-alt]{sec} والے روپ نالوں تھوڑا ودھ ابتدائی تے تھوڑا
  ودھ تفصیلی بناؤن دا مقصد اے۔
\end{editorial}

کجھ سیٹ، جیویں مثبت صحیح عدداں دا سیٹ $\PosInt$، لامتناہی نیں۔
ہُن تک اسی لامتناہی سیٹاں دیاں جیہڑیاں مثالاں ویکھیاں، اوہ سارے
!!{enumerable} سن۔ پر کجھ لامتناہی سیٹ ایہ خاصیت نہیں رکھدے۔ ایہو
جہے سیٹاں نوں \emph{!!{nonenumerable}} آکھیا جاندا اے۔

سب توں پہلاں، شاید ایہ گل پہلے ای حیران کرن والی اے کہ
!!{nonenumerable} سیٹ موجود نیں۔ ہر !!{enumerable} سیٹ~$A$ لئی کوئی
!!a{surjective} فنکشن $f \colon \PosInt \to A$ ہوندا اے۔ جے کوئی سیٹ
!!{nonenumerable} اے، تاں ایہو جہا کوئی فنکشن نہیں ہوندا۔ یعنی
$\PosInt$ دے لامتناہی !!{element}s نوں~$A$ ول نقشہ بند کرن والا کوئی
وی فنکشن سارے~$A$ نوں پورا نہیں کر سکدا۔ سو~$A$ دے !!{element}s،
لامتناہی مثبت صحیح عدداں نالوں وی ”ودھ“ نیں۔

اسی ایہ کِویں ثابت کراں گے کہ کوئی سیٹ !!{nonenumerable} اے؟ تہانوں
وکھاؤنا پئے گا کہ ایہو جہا کوئی ہر ہدف قدر لین والا فنکشن موجود نہیں
ہو سکدا۔ ایس دے برابر، تہانوں وکھاؤنا پئے گا کہ~$A$ دے عنصراں نوں اک
پاسے لامتناہی فہرست وچ گنتی وار نہیں رکھیا جا سکدا۔ ایہ کرن دا سب
توں چنگا طریقہ ایہ وکھاؤنا اے کہ~$A$ دے !!{element}s دی ہر فہرست
گھٹ توں گھٹ اک عنصر ضرور چھڈ دیندی اے؛ یا کوئی فنکشن
$f\colon \PosInt \to A$، !!{surjective} نہیں ہو سکدا۔ اسی ایہ کم
کانتور دے \emph{قطری طریقے} نال کر سکدے آں۔~$A$ دے !!{element}s دی
کوئی فہرست، مثلاً $x_1$، $x_2$، \dots، دتی ہووے، تاں اسی~$A$ دا اک
ہور عنصر بناؤندے آں جیہڑا اپنی بناوٹ کر کے اوس فہرست وچ ہو ای نہیں سکدا۔

ساڈی پہلی مثال سیٹ~$\Bin^\omega$ اے، جیہڑا $0$ تے $1$ دیاں ساریاں
لامتناہی، بنا کھلوت دیاں لڑیاں دا سیٹ اے۔

\begin{thm}
\ollabel{thm:nonenum-bin-omega}
$\Bin^\omega$، !!{nonenumerable} اے۔
\end{thm}

\begin{proof}
تضاد لئی فرض کرو کہ $\Bin^\omega$، !!{enumerable} اے، یعنی فرض کرو
کہ اک فہرست $s_{1}$، $s_{2}$، $s_{3}$، $s_{4}$، \dots{}،
~$\Bin^\omega$ دے سارے !!{element}s دی اے۔ ایہناں وچوں ہر $s_i$ آپ $0$ تے~$1$
دی اک لامتناہی لڑی اے۔ آؤ $j$ویں عنصر نوں، جیہڑا ایس فہرست دی
$i$ویں لڑی وچ اے، $s_i(j)$ آکھیے۔ تاں $i$ویں لڑی~$s_i$ ایہ اے:
\[
s_i(1), s_i(2), s_i(3), \dots
\]

اسی ایس فہرست نوں، تے ایس وچ ہر لڑی $s_i$ دے عنصراں نوں، اک جدول وچ
سجا سکدے آں:
\[
\begin{array}{c|c|c|c|c|c}
& 1 & 2 & 3 & 4 & \dots \\\hline
1 & \mathbf{s_{1}(1)} & s_{1}(2) & s_{1}(3) & s_1(4) & \dots \\\hline
2 & s_{2}(1)& \mathbf{s_{2}(2)} & s_2(3) & s_2(4) & \dots \\\hline
3 & s_{3}(1)& s_{3}(2) & \mathbf{s_3(3)} & s_3(4) & \dots \\\hline
4 & s_{4}(1)& s_{4}(2) & s_4(3) & \mathbf{s_4(4)} & \dots \\\hline
\vdots & \vdots & \vdots & \vdots & \vdots & \mathbf{\ddots}
\end{array}
\]
کھبے پاسے تھلے ول لیبل فہرست دیاں لڑیاں $s_1$، $s_2$، \dots دے
نمبر دسدے نیں؛ اُتے آر پار دتے نمبر ہر وکھری لڑی دے !!{element}s نوں
لیبل کردے نیں۔ مثال لئی، $s_{1}(1)$ اوس عدد دا ناں اے جیہڑا $0$ یا
~$1$ اے تے لڑی $s_{1}$ دا پہلا !!{element} اے؛ تے ایویں اگے۔

ہُن اسی اک لامتناہی لڑی $\overline{s}$، جیہڑی $0$ تے $1$ دی اے، بناؤندے آں،
جیہڑی ایس فہرست وچ ہو ای نہیں سکدی۔ $\overline{s}$ دی تعریف فہرست
$s_1$، $s_2$، \dots اُتے منحصر ہووے گی۔ $0$ تے $1$ دیاں لامتناہی
لڑیاں دی ہر لامتناہی فہرست اک ایہو جہی لامتناہی لڑی~$\overline{s}$
پیدا کردی اے جیہڑی پکی گل اے کہ فہرست وچ نہیں آوے گی۔

$\overline{s}$ نوں تعریف کرن لئی اسی اوہدے سارے !!{element}s دسدے آں،
یعنی اسی $\overline{s}(n)$ دسدے آں، ہر $n \in \PosInt$ لئی۔ اسی اُتلے
جدول دے قطر نوں اُتوں تھلے پڑھ کے (ایسے لئی ناں ”قطری طریقہ“ اے)، تے
فیر ہر $1$ نوں $0$ تے ہر $0$ نوں~$1$ وچ بدل کے ایہ کم کردے آں۔ ہور
مجرد طور تے، اسی $\overline{s}(n)$ نوں $0$ یا $1$ ایس مطابق تعریف
کردے آں کہ قطر دا $n$واں !!{element}، $s_n(n)$، $1$ اے یا $0$:
\[
\overline{s}(n) =
\begin{cases}
1 & \text{جے $s_{n}(n) = 0$ اے}\\
0 & \text{جے $s_{n}(n) = 1$ اے}.
\end{cases}
\]
جے تہانوں حالت وار تعریف نالوں فارمولے چنگے لگدے نیں، تاں تسیں
$\overline{s}(n) = 1 - s_n(n)$ وی تعریف کر سکدے او۔

صاف اے کہ $\overline{s}$، $0$ تے $1$ دی اک لامتناہی لڑی اے، کیوں جے
ایہ ساڈے جدول دے قطر اُتے آون والے $0$ تے $1$ دی لڑی دا بس اُلٹ عکس
اے۔ سو $\overline{s}$،~$\Bin^\omega$ دا !!a{element} اے۔ پر ایہ
فہرست $s_1$، $s_2$، \dots{} وچ نہیں ہو سکدی۔ کیوں؟

ایہ فہرست دی پہلی لڑی $s_1$ نہیں ہو سکدی، کیوں جے ایہ پہلے
!!{element} وچ $s_1$ نالوں وکھ اے۔ $s_1(1)$ جو وی ہووے، اسی
$\overline{s}(1)$ نوں اوہدے اُلٹ تعریف کیتا اے۔ ایہ فہرست دی دوجی
لڑی نہیں ہو سکدی، کیوں جے $\overline{s}$ دوجے عنصر وچ $s_2$ نالوں
وکھ اے: جے $s_2(2)$، $0$ اے، تاں $\overline{s}(2)$، $1$ اے، تے
اُلٹ صورت وچ وی ایویں۔ تے ایس طرح اگے۔

ہور ٹھیک طور تے: جے $\overline{s}$ فہرست وچ ہوندی، تاں کوئی $k$ ایہو
جہا ہوندا کہ $\overline{s} = s_{k}$۔ دو لڑیاں عین اکّو عین اوس ویلے
نیں جدوں اوہ ہر مقام اُتے اکّو قدر رکھن، یعنی ہر~$n$ لئی
$\overline{s}(n) = s_{k}(n)$۔ سو خاص طور تے $n = k$ والی حالت لیئیے،
تاں $\overline{s}(k) = s_{k}(k)$ ہونا لازمی سی۔ $s_k(k)$ یا $0$ اے
یا~$1$۔ جے ایہ $0$ اے، تاں $\overline{s}(k)$ لازماً~$1$ اے---اسی
$\overline{s}$ نوں ایویں ای تعریف کیتا اے۔ پر جے $s_k(k) = 1$ اے،
تاں فیر، $\overline{s}$ دی ساڈی تعریف کر کے، $\overline{s}(k) = 0$
اے۔ دوہاں صورتاں وچ $\overline{s}(k) \neq s_{k}(k)$۔

اسی ایہ فرض کر کے شروع کیتا سی کہ~$\Bin^\omega$ دے !!{element}s دی
اک فہرست $s_1$، $s_2$، \dots{} اے۔ ایس فہرست توں اسی اک لڑی
~$\overline{s}$ بنائی جیہڑی، جیویں اسی ثابت کیتا، فہرست وچ نہیں ہو
سکدی۔ پر ایہ وی پکی گل اے کہ اوہ $0$ تے $1$ دی اک لڑی اے جے سارے
$s_i$، $0$ تے $1$ دیاں لڑیاں نیں، یعنی
$\overline{s} \in \Bin^\omega$۔ ایہ خاص طور تے وکھاؤندا اے کہ
~$\Bin^\omega$ دے \emph{سارے} !!{element}s دی کوئی فہرست نہیں ہو
سکدی، کیوں جے ہر ایہو جہی فہرست توں اسی اک لڑی~$\overline{s}$ وی بنا
سکدے آں جیہڑی پکی گل اے کہ اوس فہرست وچ نہیں؛ سو $\Bin^\omega$ دیاں
ساریاں لڑیاں دی فہرست ہون دا مفروضہ تضاد ول لے جاندا اے۔
\end{proof}

\begin{explain}
ایس ثبوت دے طریقے نوں ”قطری بناوٹ“ آکھیا جاندا اے، کیوں جے ایہ
~$\overline{s}$ نوں تعریف کرن لئی جدول دا قطر ورتدا اے۔ قطری بناوٹ
وچ جدول دا موجود ہونا لازمی نہیں: اسی اوسے ورگا خیال ورت کے سیٹاں دا
!!{enumerable} نہ ہونا وی وکھا سکدے آں، بھاویں کوئی جدول تے حقیقی قطر
شامل نہ ہووے۔
\end{explain}

\begin{thm}
\ollabel{thm:nonenum-pownat}
$\Pow{\PosInt}$، !!{enumerable} نہیں اے۔
\end{thm}

\begin{proof}
اسی اوسے ڈھنگ نال کم کردے آں:~$\PosInt$ دے ذیلی سیٹاں دی ہر
فہرست لئی وکھاؤندے آں کہ~$\PosInt$ دا کوئی ایہو جہا ذیلی سیٹ موجود اے
جیہڑا فہرست وچ نہیں ہو سکدا۔ فرض کرو کہ~$\PosInt$ دے ذیلی سیٹاں دی
اگلی فہرست دتی ہوئی اے:
\[
Z_{1}, Z_{2}, Z_{3}, \dots
\]
ہُن اسی سیٹ $\overline{Z}$ ایویں تعریف کردے آں کہ ہر
$n \in \PosInt$ لئی، $n \in \overline{Z}$ عین اوس ویلے جدوں
$n \notin Z_{n}$:
\[
\overline{Z} = \Setabs{n \in \PosInt}{n \notin Z_n}
\]
صاف اے کہ $\overline{Z}$ مثبت صحیح عدداں دا سیٹ اے، کیوں جے مفروضے
مطابق ہر~$Z_n$ وی ایہو جہا سیٹ اے، تے ایس لئی
$\overline{Z} \in \Pow{\PosInt}$۔ پر $\overline{Z}$ فہرست وچ نہیں
ہو سکدا۔ ایہ وکھاؤن لئی اسی ثابت کراں گے کہ ہر
$k \in \PosInt$ لئی، $\overline{Z} \neq Z_k$۔

سو $k \in \PosInt$ کوئی وی ہووے۔ اسی $\overline{Z}$ نوں ایویں تعریف
کیتا اے کہ ہر $n \in \PosInt$ لئی، $n \in \overline{Z}$ عین اوس
ویلے جدوں $n \notin Z_n$۔ خاص طور تے $n=k$ رکھ کے،
$k \in \overline{Z}$ عین اوس ویلے جدوں $k \notin Z_k$۔ پر ایہ
وکھاؤندا اے کہ $\overline{Z} \neq Z_k$، کیوں جے $k$ اک دا
!!a{element} اے پر دوجے دا نہیں، تے ایس لئی $\overline{Z}$ تے $Z_k$
دے !!{element}s وکھ نیں۔ کیوں جے $k$ کوئی وی سی، ایس لئی
$\overline{Z}$ فہرست $Z_1$، $Z_2$، \dots وچ نہیں۔
\end{proof}

\begin{explain}
پچھلے ثبوت وچ قطر دا ذکر نہیں آیا، پر تسیں اینوں قطری ثبوت سمجھ سکدے
او جے ایس دی تصویر ایویں بناؤ: تصور کرو کہ سیٹ $Z_1$، $Z_2$، \dots
اک جدول وچ لکھے نیں، جتھے ہر !!{element}~$j \in Z_i$ نوں~$j$ویں کالم
وچ لکھیا گیا اے۔ من لو کہ اوس فہرست دے پہلے چار سیٹ
$\{1,2,3,\dots\}$، $\{2, 4, 6, \dots\}$، $\{1,2,5\}$، تے
$\{3,4,5,\dots\}$ نیں۔ تاں جدول ایویں شروع ہووے گا:
\[
\begin{array}{r@{}rrrrrrr}
  Z_1 = \{ & \mathbf{1}, & 2, & 3, & 4, & 5, & 6, & \dots\}\\
  Z_2 = \{ &  & \mathbf{2}, &  & 4, &  & 6, & \dots\}\\
  Z_3 = \{ & 1, & 2, &  &  & 5\phantom{,} &  & \}\\
  Z_4 = \{ &  &  & 3, & \mathbf{4}, & 5, & 6, & \dots\}\\
  \vdots & & & & & \ddots
\end{array}
\]
فیر $\overline{Z}$ اوہ سیٹ اے جیہڑا قطر دے نال تھلے جا کے، قطر اُتے
آون والے عدد چھڈ کے، تے اوہ $j$ شامل کر کے ملدا اے جتھے جدول دی
$j$ویں قطار/کالم وچ کھلوت اے۔ اُتلی صورت وچ اسی $1$ تے $2$ چھڈ
دیاں گے،~$3$ شامل کراں گے،~$4$ چھڈ دیاں گے، تے ایویں اگے۔
\end{explain}

\begin{prob}
قطری دلیل نال وکھاؤ کہ $\Pow{\Nat}$، !!{nonenumerable} اے۔
\end{prob}

\begin{prob}\label{sfr:siz:nen:prob:f-posint}
کھلی قطری دلیل نال وکھاؤ کہ فنکشناں $f \colon \PosInt \to \PosInt$
دا سیٹ !!{nonenumerable} اے۔ یعنی وکھاؤ کہ جے $f_1$، $f_2$، \dots
فنکشناں دی فہرست اے تے ہر $f_i\colon \PosInt \to \PosInt$، تاں کوئی
$\overline{f}\colon \PosInt \to \PosInt$ ایہو جہا اے جیہڑا ایس فہرست
وچ نہیں۔
\end{prob}

\end{document}
